\documentclass{report}
\usepackage{amsmath,amssymb}
\setlength{\parindent}{0mm}
\setlength{\parskip}{1em}
\begin{document}
\begin{center}
\rule{6in}{1pt} \
{\large Luz Angelica Caudillo Mata \\
{\bf A Multiscale Finite Volume Method for Geophysical Electromagnetic Simulations in Time Domain}}

4013-2207 Main Mall \\ Vancouver B C  \\ V6T 1Z4
\\
{\tt lcaudill@eos.ubc.ca}\\
Eldad Haber\end{center}

Efficient and accurate simulation of transient electromagnetic fields and
fluxes of complex geological settings is crucial in a variety of
scenarios, including mineral and petroleum exploration, water resource
utilizations and geothermal power extractions. In this talk, we present a
new Multiscale finite volume method that computes these electromagnetic
responses in a very effective manner.

Geophysical time-varying electromagnetic simulations of highly
heterogeneous media are computationally expensive. One reason for this is
the fact that very fine meshes are often required to accurately
discretize the physical properties of the media, which may vary over a
wide range of length scales and several orders of magnitude. Using very
fine meshes for the discrete models lead to solve large systems of
equations that are often difficult to deal with at each time step.

To reduce the computational cost of the electromagnetic simulation, we
develop a Multiscale finite volume method for the quasistatic Maxwell�s
equations in the time domain. Our method begins by locally computing
Multiscale basis functions at each time step, which incorporate the
small-scale information contained in the physical properties of the
media. Using a Galerkin proper orthogonal decomposition approach, the
local basis functions are used to represent the solution on a coarse
mesh. The governing equations are numerically integrated using an
implicit time marching scheme.

Our approach leads to a significant reduction in the size of the final
system of equations to be solved and in the amount of computational time
of the simulation, while accurately approximating the behavior of the
fine-mesh solutions. We demonstrate the performance of our method in the
context of an airborne geophysical application.


\end{document}
